\section{第一波开放 Instruct 模型：数据分布比名字更重要}

上一章定义了 IFT，本章跟随 2023 年第一波模型。Alpaca 证明合成 instruction data 可以快速适配 LLaMA；Vicuna 证明真实聊天 prompt 分布很关键；OpenAssistant 证明开放人类数据有长期价值；StableVicuna 说明 PPO 并非只有闭源公司能运行。核心变量不断从“有没有算法”转向“有没有对的数据”。

\subsection{Alpaca 与 instruction fine-tuning}

本节从 Alpaca 开始。它的重要性不在于今天的绝对质量，而在于展示：一个强 base model 加上规模不大的 instruction-response data，就能获得明显 chat/instruct 行为。读图时先区分 base model 的知识与 fine-tuning 学到的响应格式。

\lecturefigure{slide-029.jpg}{Alpaca：第一批广泛传播的开放 instruction-tuned model}{official deck, p.~29}

\lecturefigure{slide-030.jpg}{IFT 的两个目标：适配输入风格并支持 system prompt 与 multi-turn}{official deck, p.~30}

若只在 assistant answer token 上计算损失，IFT 目标可写成
\[
\mathcal{L}_{\mathrm{IFT}}=-\sum_{t\in\mathcal{A}}\log p_{\theta}(y_t\mid s,u,y_{<t}),
\]
其中 $s$ 是 system prompt，$u$ 是 user instruction，$y$ 是 assistant response，$\mathcal{A}$ 是需要监督的 assistant token 位置。Chat template 把角色和边界编码成特殊 token；模板错配会让同一模型表现显著下降。

\lecturefigure{slide-031.jpg}{IFT 从 base model 出发，用 instruction-response pairs 继续训练}{official deck, p.~31}

\begin{importantbox}{IFT 主要改变什么}
IFT 通常强化三类能力：识别“这是一个请求”、按照角色/格式组织回答、把已有知识映射到指令需要的输出。它可能改善任务能力，也可能遗忘或偏移 base capability；因此需要同时评估 chat behavior 与原始能力，而不能把“更会聊天”当作“所有方面更强”。
\end{importantbox}

\teachervoice{00:13:42--00:14:53，Nathan 把 IFT 描述成让模型整合 chat template、system prompt、多轮角色与输入风格。它仍使用 autoregressive loss，差别主要来自训练样本的条件结构。}

\subsection{Self-instruct：用模型扩张 instruction data}

上一节说明 IFT 的行为变化主要来自 instruction-response 分布；本节进一步追问：开放团队没有大规模人工标注预算时，这个分布从哪里来。Alpaca 的可复制性来自 self-instruct：先用少量人工 seed prompt，让强模型生成更多指令与回答，再过滤重复、低质量或格式错误样本。读图时要把 seed 质量、生成多样性和过滤规则看成同一个数据系统；它把昂贵人工编写变成“人工种子 + 模型扩增 + 过滤”的 pipeline，也带来教师模型偏见和分布收缩风险。

\lecturefigure{slide-032.jpg}{Self-instruct：从少量高质量 seed 扩张合成 instruction data}{official deck, p.~32}

设 seed 集合为 $D_0$，生成器为 $G$，过滤器为 $F$，则一轮扩张可写成
\[
D_{k+1}=D_k\cup F\bigl(G(D_k)\bigr).
\]
其中 $G(D_k)$ 生成新 instruction/response，$F$ 去重并执行质量规则。若 $G$ 和 $F$ 只偏好相似模板，数据规模增加也不会增加真正的任务覆盖。

\begin{warningbox}{Synthetic data 的两个闭环风险}
第一，student 会继承 teacher 的事实错误、拒答边界和语言风格；第二，反复 self-generation 会压缩多样性，使模型越来越擅长“模型式问题”。因此应混入真实用户和人类数据，并测量 prompt/response 的语义覆盖，而不只统计样本数。
\end{warningbox}

\teachervoice{00:14:53--00:16:10，课程把 self-instruct 视为早期开放模型爆发的关键：生成回答的模型替代了部分人工标注，使 instruction data 变得可获得。后半场又提醒合成分布可能重复而脆弱。}

\subsection{Vicuna 与 ShareGPT：真实用户 prompt 改变分布}

Self-instruct 解决“怎样快速扩量”，但还没有回答“训练问题是否像真实用户会问的问题”。Vicuna 因而把关注点从样本数量转向 prompt source：ShareGPT 收集用户分享的 ChatGPT 对话，使开放模型第一次看到更接近真实产品的多轮请求、追问和风格。看接下来的两页时，应同时比较 synthetic instruction 与真实会话在任务长度、轮次结构和隐私风险上的差异；质量提升来自分布匹配，而不只是更多 token。

\lecturefigure{slide-033.jpg}{Vicuna 在 Alpaca 之后迅速加入更真实的 chat prompt distribution}{official deck, p.~33}

\lecturefigure{slide-034.jpg}{ShareGPT 数据来源、价值与隐私风险}{official deck, p.~34}

\begin{knowledgebox}{读图：数据价值与数据治理同时出现}
ShareGPT 提供了真实用户想问什么，这是 synthetic instruction 难以复制的信号；但共享工具并不天然等于用户同意用于训练，文本还可能包含个人信息。后来的 LMSYS-Chat-1M 与 WildChat 更重视清洗和 consent，说明开放数据的合法性与质量是同一 pipeline 的要求。
\end{knowledgebox}

\teachervoice{00:17:21--00:18:55，Nathan 说 ShareGPT 是少数能让开放 builder 看到真实用户 prompt 的来源，并指出后来的数据集花了大量工作清理个人信息或在收集前取得同意。}

\subsection{Weight differences 与开放发布约束}

前两节讨论的是数据如何获得，本节转向另一个常被忽略的系统约束：训练完成并不等于模型可以合法、稳定地交付。LLaMA 最初的权重许可和传播方式使 fine-tune 发布变得尴尬，团队只能发布相对 base model 的 weight difference，用户必须合法获得 base weights 再合并。读这组时间线时要区分模型能力、训练配方与分发机制；weight diff 不是机器学习算法，却直接决定模型是否可安装、可复现和可迭代。

\lecturefigure{slide-035.jpg}{开放 instruct 模型继续扩张：Vicuna、Koala、Dolly 等}{official deck, p.~35}

\lecturefigure{slide-036.jpg}{为什么发布 weight differences：在许可约束下分发 fine-tune}{official deck, p.~36}

若 base 权重为 $W_0$、fine-tuned 权重为 $W_1$，发布差分就是
\[
\Delta W=W_1-W_0,\qquad W_1=W_0+\Delta W.
\]
其中 $\Delta W$ 并不减小完整 fine-tune 的信息量，只把分发依赖转移给用户。它与 LoRA 的低秩 update 不同：weight diff 可以是全秩、同尺寸张量。

\lecturefigure{slide-037.jpg}{模型时间线继续加入 Koala 与 Dolly，发布方式逐步标准化}{official deck, p.~37}

\lecturefigure{slide-038.jpg}{早期开放 instruct 模型的 MT-Bench 分数变化}{official deck, p.~38}

\begin{warningbox}{历史分数只能解释当时生态}
MT-Bench 的 judge、prompt 和模型版本后来都可能更新。图中分数用于说明 2023 年模型相对变化，不能直接与 2026 年模型或不同 judge 配置横比。复现时应记录评测 commit、judge model、temperature 和 prompt template。
\end{warningbox}

\teachervoice{00:19:00--00:21:05，讲者解释 weight diff 是对 LLaMA 发布约束的工程回应，后来许可变化才让直接发布变得容易。开放生态的速度受法律与分发接口影响，不只是训练代码。}

\subsection{OpenAssistant 与 StableVicuna：人类数据和早期 PPO}

本节看两条常被简化的线索。OpenAssistant 发布人类生成、标注的对话数据，并长期成为后续模型的原料；StableVicuna 则较早用 PPO 完成开放 RLHF。二者说明开放社区并非只有合成 IFT，也在探索高成本的人类偏好和 RL pipeline。

\lecturefigure{slide-039.jpg}{OpenAssistant Conversations：早期开放人类 instruction 数据集}{official deck, p.~39}

\lecturefigure{slide-040.jpg}{StableVicuna：较早使用 PPO 的开放 RLHF 模型}{official deck, p.~40}

\begin{importantbox}{Human data 为什么难以替代}
人类数据能提供真实意图、多样语言、边界案例和相对偏好；这些信号决定模型在“合理但不唯一”的回答之间如何排序。成本来自招募、培训、一致性、隐私与质量控制，因此开放社区最稀缺的往往不是训练脚本，而是可合法发布的高质量人类数据。
\end{importantbox}

\teachervoice{00:23:43--00:24:47，Nathan 称 OpenAssistant 是早期最成功的开放项目之一，并把“需要更多 open human data”作为全讲结论；同时提醒 StableVicuna 在 2023 年就让 PPO 工作，只是能做到的团队很少。}

\subsection{本章小结}

第一波开放 instruct 模型的决定性变量依次是：可用 base weights、合成 instruction pipeline、真实聊天 prompt、人类开放数据、发布许可和可运行的 RLHF。模型名称不断更换，但数据分布和可复现接口是更稳定的解释。下一章将看 QLoRA 如何进一步降低参与门槛，以及为什么低内存 fine-tuning 不自动等于低成本 RLHF。

